\documentclass[9pt,xcolor=table]{beamer}
%add handout for no pauses
\usepackage[english]{babel}

\usepackage{amsmath}
\usepackage[utf8x]{inputenc}
\usepackage[T1]{fontenc}% útil para que ponga letras acentuadas y no junte una vocal con un acento.
\usepackage{pgfpages}
\usepackage{../style/arrowsnew}
\usepackage{ulem}
\usepackage{tikz,pgf}
\usetikzlibrary{shapes.callouts,decorations.pathmorphing,arrows,snakes,shapes,patterns}
\usepgflibrary{shapes.callouts}
\usepackage{xcolor}
\usepackage{babel}
\usepackage{tikz}
\usepackage{multicol}
\usepackage{graphicx}
\usepackage{../style/beamerthemeMPII}
\usepackage{minted}
%[xleftmargin=20pt,linenos] usual option 


\input{../style/CommonPackages}

\input{../style/BoxesEtAl}

\input{../style/NotationShortcuts}

\newcommand{\argmin}{\operatorname{argmin}}
\newcommand{\argmax}{\operatorname{argmax}}
\newcommand{\sm}{\operatorname{sm}}
%%% SOME OPTIONS

%% for title page
\newif\ifshowtitledate
% TODO teaser image

%% for status bar
\newif\ifshowshorttitle
\newif\ifshowshortauthor
\newif\ifshowdate
\newif\ifshowtotalnumberofframes

%% for content
\newif\ifshowminiframes
\newif\ifincludetoc

% default values
\showtitledatefalse
\showshorttitlefalse
\showshortauthorfalse
\showdatefalse
\showtotalnumberofframesfalse
\showminiframesfalse
\includetocfalse

%%%% TITLE
\makeatletter
\newcommand\changemode[1]{%
  \gdef\beamer@currentmode{#1}}
\makeatother
% Title page containing title, author name, date
\mode<presentation>
%%%% EDIT HERE %%%%%%%%%%%%%%%%%%%%%%%%%%%%%%%%%%%%%
% Specify the title, put the shorttitle in square brackets, the full title in curly brackets
    \title[]{DOCE-102 Statistical Computing\\
    Unit 8: Neural Networks}

    \showshorttitletrue
   \author[N.R.]{Nicol\'as Rivera}

   % \showshortauthortrue
    \date{1S 2025}
    \showtitledatetrue
     \showtotalnumberofframestrue
    \MPIIcolorscheme{MPIIdarkblue}
%%%%%%%%%%%%%%%%%%%%%%%%%%%%%%%%%%%%%%%%%%%%%%%%%%%%%%%%%%%%%%%%%%%%%%%%%%%%%%

\begin{document}


\begin{frame}
\titlepage
\begin{tikzpicture}[overlay,scale=0.45,opacity=1.0,remember picture]
%\node () at (12,6) {
%\includegraphics[scale=0.55]
%{geometry2title}};
\node () at (12,-7) {
\includegraphics[scale=0.04]
{../style/uvalparaiso}};
\end{tikzpicture}
\end{frame}
\setbeamertemplate{headline}[MPIIheadline]
\setbeamertemplate{footline}[MPIIfootline]

  \setbeamertemplate{mini frames}[tick]
  \setbeamersize{mini frame size=0.135cm}

\AtBeginSection[] % Do nothing for \section*
{
\begin{frame}
\frametitle{Table of Contents}
\tableofcontents[currentsection]
\end{frame}
}

\section{Warning Message}

\begin{frame}{Warning Message}
Neural networks are very popular, and thus the internet is full of non-useful information
\vspace{1cm}

As an introduction to the subject, I recommend Chapter 11 of the Book \vspace{0.5cm}
\begin{center}

\textit{\textcolor{red}{Elements of Statistical learning}}\\
 by Hastie, Tibshirani, and Friedman
\end{center}
\end{frame}

\section{Introduction}
\begin{frame}{A Neuron = A regression}
The basic element of a Neural Network is a Neuron. 

A neuron is, more or less, a linear regression or a logistic regression.

\fancythm{Linear Regression Model}
{
	In linear regression we model a response $Y \in \R$ in terms of covariates $X \in \R^d$ by using weights $\alpha \in \R^d$, i.e.
	$$Y = \alpha_0 + \alpha^{\transpose} X.$$
		As usual, it is convenient to include the intercept in $X$ to avoid writing the $\alpha_0$ term.
}

\fancythm{logistic Regression Model}
{
	In logistic regression we model the probability of a response $Y \in \{0,1\}$ in terms of covariates $X \in \R^d$ by using weights $\alpha \in \R^d$, and a link function $\sigma:\R\to [0,1]$
	$$\Prob(Y=1) = \sigma(\alpha_0 + \alpha^{\transpose} X).$$
 Usually, $\sigma(x) = \frac{1}{1+e^{-x}}.$
}
\end{frame}

\begin{frame}{A Neuron = A regression}
Sometimes is useful to extend logistic regression to $K$ classes by using the so-called soft-max function $\sm:\R^K \to \R^K$ defined by

$$\sm_k(T) = \frac{e^{T_k}}{\sum_{\ell=1}^K e^{T_\ell}}.$$

In this case, we can extend logistic regression by $K$ groups.

\fancythm{K-groups extension of logistic regression}
{
	Consider $X \in \R^d$, define 
	
	 For $k \in \{1,\ldots, K\}$, let $\alpha_k \in \R^d$. Then K-groups extension of logistic regression is defined by
	 
	 \begin{align*}
	 \Prob(Y=k) &= \sm_k(\alpha_{k0} + \alpha_k^{\transpose} X),  \text{ for $k \in \{1,\ldots, K\}$}
	 \end{align*}
	 
	 %Usually, to avoid identifiability issues, we  can choose $\alpha_1$ as the vector with only 1's. 
}
\textbf{Exercise:} verify that when $K=2$ we recover exactly the logistic regression with link $\sigma(x) = 1/(1+e^{-x})$.
\end{frame}

\begin{frame}{1-hidden-layer Neural Networks}

\fancythm{1-hidden-layer NN for Classification}
{
Suppose we have a data point $(X,Y)$ where $Y \in \{1,\ldots, K\}$ and $X \in \R^d$. We can write the following model for $\Prob(Y)$ in terms of $X$.

\begin{align*}
Z_m &= \sigma(\alpha_{m0} + \alpha_i^{\transpose} X), \text{ for } m \in \{1,\ldots, M\}\\
\Prob(Y=k) &= \sm_k(\beta_{k0} +  \beta_k^{\transpose} Z), \text{ for } k \in \{1,\ldots, K\}
\end{align*}
where $\alpha_{m} \in \R^d$ and $\beta_k \in \R^M$.
}
\vspace{1cm}

A few comments:

\begin{enumerate}
\item Here the $Z_m$ is so-called hidden layer, because we do not observe them.
\item In practice, a few parameters can be fixed to make the model identifiable. For example, we can set $\beta_1$ as the vector with only 1's. In practice practice, this solution is an extra problem.
\item We denote by $\mathbf{\alpha}$ the matrix containing the $\alpha$'s, each $\alpha_m$ in a row (including $\alpha_{m0}$)
\item We denote by $\mathbf{\beta}$ the matrix containing the $\beta$'s, each $\beta_k$ in a row (including the intercept)
\end{enumerate}
\end{frame}
\begin{frame}{1-hidden-layer Neural Networks}

\begin{center}
\includegraphics[scale=0.8]{NN_classification.pdf} 
\end{center}
\end{frame}

\begin{frame}{1-hidden-layer Neural Networks}

\begin{center}
\includegraphics[scale=0.6]{NN_classification.pdf} 
\end{center}
\pause
\begin{enumerate}
\item The middle layer is called the hidden layer because we do not observe the variables $Z$
\item Each $Z_m$ is called a (hidden) neuron
\item $\alpha$ and $\beta$ are called the weight matrices
\end{enumerate}
\end{frame}

\begin{frame}{1-hidden-layer Neural Networks}
We can also have a model with a continuous response $Y$, with a very similar picture

\fancythm{1-hidden-layer NN for Regression}
{
Suppose we have a data point $(X,Y)$ where $Y \in \R^K$ and $X \in \R^d$. Then

\begin{align*}
Z_m &= \sigma(\alpha_{m0} + \alpha_i^{\transpose} X), \text{ for } m \in \{1,\ldots, M\}\\
\textcolor{violet}{Y_k }&= \textcolor{violet}{\beta_{k0} +  \beta_k^{\transpose} Z, \text{ for } k \in \{1,\ldots, K\}}
\end{align*}
where $\alpha_{m} \in \R^d$ and $\beta_k \in \R^M$.
}
\begin{center}
\includegraphics[scale=0.47]{NN_regression.pdf} 
\end{center}
\end{frame}

\begin{frame}{1-hidden-layer Neural Networks}
The good thing about neural networks is that they are very flexible. \\

Consider a 1-hidden-layer NN for regression. Denote by
$\alpha$ and $\beta$ its weight matrices, and by $f_{{\alpha},{\beta}}(x)$ the output when the input is $x$, i.e. $f_{{\alpha},{\beta}}$ is a function $\R^d\to \R^K$.\\

The following approximation results is due to Cybenko\footnote{Cybenko G -  Approximation by Superposition of a Sigmoidal Function}


\begin{theorem}

\textbf{\textit{Informal Version:}} Every continuous-function $f:[0,1]^d \to \R^K$ can be approximated by a 1-hidden-layer Neural Network for regression.\\

\vspace{1cm}

\textbf{\textit{Formal Version:}} For every continuous-function $f:[0,1]^d \to \R^K$ and every $\varepsilon>0$, it exists a NN network $f_{\alpha,\beta}$ such that

$$\|f(x)-f_{\alpha,\beta}(x)\|_2< \varepsilon$$
for all $x \in \R^d$.
\end{theorem}
\end{frame}

\begin{frame}{Fitting a Neural Network}

In a 1-hidden-layer Neural Network for covariate $X \in \R^d$ and response $Y \in \R^K$ we have the parameters

$$\alpha = (\alpha_{m \ell}: \ell \in \{0,\ldots, d\}, m \in \{1,\ldots, M\})$$
and
$$\beta = (\beta_{km}: m \in \{0,\ldots, M\}, k \in \{1,\ldots, K\})$$

so there are $M(d+1)+K(M+1)$ parameters.

\vspace{1cm}

Denote by $f_{\alpha,\beta}(x)$ the output of such network when the input is $x\in \R^d$.

\end{frame}

\begin{frame}{Fitting a Neural Network}
Suppose we now observe data $(X_i,Y_i)_{i=1}^n$. We want to fit the best parameters. 

\begin{enumerate}
\item For regression, we use the typical loss function
$$L(\alpha,\beta) = \sum_{i=1}^n \sum_{k=1}^K (Y_{ik}-f_{k, \alpha,\beta}(X_i))^2$$
\item For classification we can use 
$$L(\alpha,\beta) = \sum_{i=1}^n \sum_{k=1}^K (Y_{ik}-f_{k,\alpha,\beta}(X_i))^2$$
or the so-called cross-entropy
$$L(\alpha,\beta) = -\sum_{i=1}^n \sum_{k=1}^K Y_{ik} \log f_{k,\alpha,\beta}(X_i)$$
\end{enumerate}

\pause

In general, we may want to add some penalisation term, otherwise a NN will overfit our data.
\end{frame}



\section{Gradient of a NN}


\begin{frame}{Gradient of a NN}

By being very careful, we can compute derivatives associated with the loss function of the NN. A general approach is to use the back propagation algorithm (see extra material)

\vspace{0.5cm}

\begin{center}
\includegraphics[scale=0.5]{NN_classification.pdf} 
\end{center}
\end{frame}

\begin{frame}{Fitting a Neural Network}
Recall that we observe data $(X_i,Y_i)_{i=1}^n$, and we use loss functions

\begin{enumerate}
\item For regression, we use the typical loss function
$$L(\alpha,\beta) = \sum_{i=1}^n \sum_{k=1}^K (Y_{ik}-f_{k, \alpha,\beta}(X_i))^2$$
\item For classification we can use 
$$L(\alpha,\beta) = \sum_{i=1}^n \sum_{k=1}^K (Y_{ik}-f_{k,\alpha,\beta}(X_i))^2$$
or the so-called cross-entropy
$$L(\alpha,\beta) = -\sum_{i=1}^n \sum_{k=1}^K Y_{ik} \log f_{k,\alpha,\beta}(X_i)$$
\end{enumerate}

Also, we might add a penalisation term.
\end{frame}

\begin{frame}
Some consideration:
\begin{enumerate}
\item Recall that towards the first layer we always use the sigmoid function $\sigma$
\item The link function between first and second layer may change depending on the problem. For a generic explanation, call it $g:\R^K \to \R^K$, with $g(t) = (g_1(t),\ldots, g_K(t))$.
\item It is convenient to add a 1 in front of each covariate vector, so we write $\alpha_m^{\transpose} X$ instead of $\alpha_{m0}+ \alpha_m^{\transpose} X$, etc.
\item for ease of notation, instead of writting $f_{k,\alpha,\beta}(x)$ for the $k$-th output of the NN when the input is $x$, we just write $f_k(x)$.
\end{enumerate}

\end{frame}

\begin{frame}{Gradient of a NN: Chain Rule}
		\fancythm{Chain Rule}{
		Let $x_1,\ldots, x_d$ be variables in $\R$. Let $a_1,\ldots, a_m:\R^d \to \R$ be functions, and let $F:\R^m \to \R$ be another function.\\
		
		Let $g(x_1,\ldots,x_d) = F(a_1(x_1,\ldots, x_d), a_2(x_1,\ldots, x_d),\ldots, a_m(x_1,\ldots, x_d))$\\
		
		\begin{align*}
			&\left(\frac{\partial g}{\partial x_k}\right)g(x_1,\ldots,x_d)\\
			= &\sum_{i=1}^r \left(\frac{\partial a_i}{\partial x_k}\right)(x_1,\ldots, x_d) \times  \left(\frac{\partial F}{\partial a_i}\right)(a_1(x_1,\ldots, x_d), a_2(x_1,\ldots, x_d),\ldots, a_m(x_1,\ldots, x_d))
		\end{align*}
	}\textbf{}
\end{frame}
\begin{frame}{Gradient of a NN}

\begin{enumerate}
\item Write $L_i = \sum_{k=1}^K (Y_{ik} - f_k(X_i))^2$, implicitly depending on $\alpha$ and $\beta$

\item Let $Z_{im} =\sigma(\alpha_m^{\transpose} X_i)$ and $Z_i = (Z_{i1},\ldots, Z_{iM})$, and recall that $X_i = (1, X_{i1},\ldots, X_{id})$
\item Let $T_{ik} = \beta_{k}^{\transpose} Z_i$ and $T_i = (T_{i1},\ldots, T_{iK})$

\item It is important to know that
\begin{align*}
\sigma'(x)& = \sigma(x)(1-\sigma(x))\\
\frac{\partial \sm_k(t)}{\partial t_{\ell}} &= \begin{cases}\sm_k(t)\sm_{\ell}(t) & \text{ if } k\neq \ell\\
-\sm_k(t)(1-\sm_k(t)) & \text{ if } k=\ell\end{cases}
\end{align*}
\item Also recall the chain rule for multivariate functions. If $r = r(a) \in \R^d$ with $a \in \R$,  and $F:\R^d \to R$, then 
$$\frac{\partial F(r(a))}{ \partial a} =  \sum_{\ell=1}^d \frac{\partial r_\ell(a)}{\partial a} \left(\frac{\partial }{ \partial r_{\ell}} F\right)(r(a))$$
\item Other useful facts: $Z_{im}$ only depends on $\alpha_{\ell m}$ for $\ell \in \{0,\ldots, d\}$, and $f_k(X_i)$ only depends on $\beta_{mk}$ for $m\in \{0,\ldots, M\}$

\end{enumerate}

\end{frame}

\begin{frame}{Gradient of a NN}
For classification, we have $g_k(T) = \sm_k(T)$. Write $L_i = \sum_{k=1}^K (Y_{ik} - f_{k}(X_i))^2$, then

\begin{align*}
\frac{\partial L_i}{\partial \alpha_{m \ell}} &= -2\sum_{k=1}^K(Y_{ik}-f_k(X_i))  \sum_{j=1}^m \frac{\partial Z_{ij}}{\partial \alpha_{ m \ell}} \frac{\partial g_k(T_i)}{\partial Z_{ij}}\\
&= -2\sum_{k=1}^K(Y_{ik}-f_k(X_i))  \frac{\partial Z_{im}}{\partial \alpha_{ m \ell}} \frac{\partial g_k(T_i)}{\partial Z_{im}}\\
\end{align*}
now, $\frac{\partial Z_{im}}{\partial \alpha_{m \ell}} = X_{i\ell} Z_{im}(1-Z_{im})$ and 
\begin{align*}
\frac{\partial g_k(T_i)}{\partial Z_{im}} &= \sum_{j=1}^K \frac{\partial T_{ij}}{\partial Z_{im}} \frac{\partial g_k(T_i)}{\partial T_i}\\
&=\sum_{j=1}^K \beta_{jm} \frac{\partial g_k(T_i)}{\partial T_{ij}} = \sum_{j=1}^K \beta_{jm} \frac{\partial \sm_k(T_i)}{\partial T_{ij}} 
\end{align*}

\end{frame}

\begin{frame}{Gradient of a NN}
Also, 

\begin{align*}
\frac{\partial L_i}{\partial \beta_{jm}} &= -2\sum_{k=1}^K (Y_{ik} - f_k(X_i)) \sum_{p=1}^K \frac{\partial T_{ip}}{\partial \beta_{jm}} \frac{\partial g_k(T_i)}{\partial T_{ip}}\\
&= -2\sum_{k=1}^K (Y_{ik} - f_k(X_i)) Z_{im} \frac{\partial \sm_k(T_i)}{\partial T_{ij}}
\end{align*}
\end{frame}

\begin{frame}{Gradient of a NN}
With this, we can do a DG algorithm!
We update as following:

\begin{align*}
\beta_{km}^{t+1} &= \beta_{km}^t- \gamma_t \sum_{i=1}^n \frac{\partial L_i^t}{\partial \beta_{km}^t}\\
\alpha_{m\ell}^{t+1} &= \alpha_{m\ell} - \gamma_t \sum_{i=1}^n \frac{\partial L_i^t}{\partial \alpha_{km}^t}
\end{align*}

%where is the Back-propagation?

%\begin{enumerate}
%\item Given initial parameters $\alpha^t$ and $\beta^t$, we have to %evaluate $R$ (that's a forward evaluation)
%\item Then we have to \textit{update} back to obtain $\alpha^{t+1}$ and $%\beta^{t+1}$.%
%\end{enumerate}
\end{frame}

\begin{frame}{Gradient of a NN}
\textbf{Exercise:} Compute the gradient of a NN for Classification when the loss function is $L(\alpha,\beta) = -\sum_{i=1}^n \sum_{k=1}^K Y_{ik} \log f_k(X_i)$\\

\vspace{1cm}

\textbf{Exercise 2:} Compute the gradient of NN for \textbf{Regression}, i.e. the loss function is $L(\alpha,\beta) = \sum_{i=1}^n \sum_{k=1}^K (Y_{ik}-f_k(X_i))^2$ and $g_k(t) = t$.
\end{frame}

\section{Material Extra}

\section{Back Propagation}
\begin{frame}
	Since NNs have so many parameters and nested functions we need to be very careful at computing derivatives and applying chain rules.
	\\
	
	\vspace{1cm}
	
	We first need a nice representation of functions. For that we consider a network of functions
	
\end{frame}

\begin{frame}{Network of functions}
	
	A network of function is 
	
	\begin{enumerate}
		\item A graph with vertices $V$ and directed edges $E$ with no directed cycles
		\item A set of initial vertices $S\subseteq V$ and each $s \in S$ has an associated value $x_s$. Vertices in $S$ only have outgoing edges
		\item Each non-initial vertex has associated a function $f:\R^d \to \R$ where $d$ is the number  of incoming vertices of the vertex
		\item A final vertex $f \in V$ that has no outgoing edges.
		\item \textbf{Requirement:} every vertex $v \in V$ belongs to a directed path from a vertex $s \in V$ to $f$.
		\item \textbf{For implementation purposes:} each vertex has a list of incoming vertices and outgoing vertices. To evaluate functions variables are introduced into the function in the order given by the incoming list.
	\end{enumerate}
	
\end{frame}

\begin{frame}{Network of functions}
	\includegraphics[scale=0.8]{FunctionNetwork.pdf} 
\end{frame}



\begin{frame}{Network of functions: Evaluation}
	The evaluation of a network of function is quite easy:
	
	\begin{itemize}
		\item Set the initial values for the vertices in $S$, and the other vertices are without value
		\item A vertex is \textbf{ready} if it has a value
		\item Every vertex that is ready pass the value to its neighbours following the outgoing edges
		\item If a vertex has received values by all the vertices in the incoming list, then it evaluates the function and set its value to the outcome of the function
	\end{itemize} 
\end{frame}


\begin{frame}{Network of functions}
	\includegraphics[scale=0.8]{FunctionNetwork1.pdf} 
\end{frame}


\begin{frame}{The back propagation algorithm}
	
	Since these network represent composition of networks they are ideal for the chain rule.
	
	\fancythm{Chain Rule}{
		Let $x_1,\ldots, x_d$ be variables in $\R$. Let $a_1,\ldots, a_m:\R^d \to \R$ be functions, and let $F:\R^m \to \R$ be another function.\\
		
		Let $g(x_1,\ldots,x_d) = F(a_1(x_1,\ldots, x_d), a_2(x_1,\ldots, x_d),\ldots, a_m(x_1,\ldots, x_d))$\\
		
		\begin{align*}
			&\left(\frac{\partial g}{\partial x_k}\right)g(x_1,\ldots,x_d)\\
			= &\sum_{i=1}^r \left(\frac{\partial a_i}{\partial x_k}\right)(x_1,\ldots, x_d) \times  \left(\frac{\partial F}{\partial a_i}\right)(a_1(x_1,\ldots, x_d), a_2(x_1,\ldots, x_d),\ldots, a_m(x_1,\ldots, x_d))
		\end{align*}
	}
	
\end{frame}


\begin{frame}{The back propagation algorithm}
	\begin{center}
		\includegraphics[scale=0.45]{FN2.pdf} 
	\end{center}
	We want the gradient of  $g_6(x,y,z) = f_6(f_4( f_1(x,y,z),f_2(x,y,z) ), f_5(f_2(x,y,z),f_3(x,y,z))$. Denote
	
	\begin{align*}
		g_5(x,y,z) &= f_5(f_2(x,y,z),f_3(x,y,z))\\
		g_4(x,y,z) &= f_4( f_1(x,y,z),f_2(x,y,z) )\\
		g_3(x,y,z) &= f_3(x,y,z)\\
		g_2(x,y,z) &= f_2(x,y,z)\\
		g_1(x,y,z) &= f_1(x,y,z)
	\end{align*}
\end{frame}

\begin{frame}{The back propagation algorithm}
	\begin{center}
		\includegraphics[scale=0.45]{FN3.pdf} 
	\end{center}
	We want the gradient of  $g_6(x,y,z) = f_6(f_4( f_1(x,y,z),f_2(x,y,z) ), f_5(f_2(x,y,z),f_3(x,y,z))$. Denote
	
	\begin{align*}
		g_5(x,y,z) &= f_5(f_2(x,y,z),f_3(x,y,z))\\
		g_4(x,y,z) &= f_4( f_1(x,y,z),f_2(x,y,z) )\\
		g_3(x,y,z) &= f_3(x,y,z)\\
		g_2(x,y,z) &= f_2(x,y,z)\\
		g_1(x,y,z) &= f_1(x,y,z)
	\end{align*}
\end{frame}

\begin{frame}
	\begin{center}
		\includegraphics[scale=0.45]{FN4.pdf} 
	\end{center}
	
	\begin{align*}
		\left(\frac{\partial g_6}{\partial x}\right)(x,y,z)& = \textcolor{violet}{\left(\partial_1 f_6 \right)}(g_4(x,y,z), g_5(x,y,z)) \times \frac{\partial g_4}{\partial x}(x,y,z) \\
		&+ \textcolor{violet}{\left(\partial_2 f_6\right)}(g_4(x,y,z), g_5(x,y,z)) \times \left(\frac{\partial g_5}{\partial x}\right)(x,y,z)
	\end{align*}
	
	Here $\partial_i f$ represents partial derivative with respect to the $i$-th coordinate of $f$.
\end{frame}


\begin{frame}
	\begin{center}
		\includegraphics[scale=0.45]{FN5.pdf} 
	\end{center}
	
	\begin{align*}
		\left(\frac{\partial g_6}{\partial x}\right)(x,y,z)& = \textcolor{violet}{\left(\partial_1 f_6 \right)}(g_4(x,y,z), g_5(x,y,z)) \times \frac{\partial g_4}{\partial x}(x,y,z) \\
		&+ \textcolor{violet}{\left(\partial_2 f_6\right)}(g_4(x,y,z), g_5(x,y,z)) \times \left(\frac{\partial g_5}{\partial x}\right)(x,y,z)\\
		&=\textcolor{red}{h_{46}} \times \frac{\partial g_4}{\partial x}(x,y,z)+ \textcolor{red}{h_{56}} \times \left(\frac{\partial g_5}{\partial x}\right)(x,y,z)
\end{align*}\end{frame}

\begin{frame}
	Iteratively, we find the derivatives of $g_4$ and $g_5$. \\
	
	For $g_4$ we have
	
	\begin{align*}
		\left(\frac{\partial g_4}{\partial x}\right)(x,y,z)& = \textcolor{violet}{\left(\partial_1 f_4 \right)}(g_1(x,y,z), g_2(x,y,z)) \times \frac{\partial g_1}{\partial x}(x,y,z) \\
		&+ \textcolor{violet}{\left(\partial_2 f_4\right)}(g_1(x,y,z), g_2(x,y,z)) \times \left(\frac{\partial g_2}{\partial x}\right)(x,y,z)\\
		&=\textcolor{red}{h_{14}} \times \frac{\partial g_1}{\partial x}(x,y,z)+ \textcolor{red}{h_{24}} \times \left(\frac{\partial g_2}{\partial x}\right)(x,y,z)
	\end{align*}
	\begin{center}
		\includegraphics[scale=0.45]{FN6.pdf} 
	\end{center}
\end{frame}

\begin{frame}
	Similarly, for $g_5$
	
	\begin{align*}
		\left(\frac{\partial g_5}{\partial x}\right)(x,y,z)& = \textcolor{violet}{\left(\partial_1 f_5 \right)}(g_2(x,y,z), g_3(x,y,z)) \times \frac{\partial g_2}{\partial x}(x,y,z) \\
		&+ \textcolor{violet}{\left(\partial_2 f_5\right)}(g_2(x,y,z), g_3(x,y,z)) \times \left(\frac{\partial g_3}{\partial x}\right)(x,y,z)\\
		&=\textcolor{red}{h_{25}} \times \frac{\partial g_2}{\partial x}(x,y,z)+ \textcolor{red}{h_{35}} \times \left(\frac{\partial g_3}{\partial x}\right)(x,y,z)
	\end{align*}
	\begin{center}
		\includegraphics[scale=0.45]{FN7.pdf} 
	\end{center}
	
	
	
	
\end{frame}

\begin{frame}
	Up to this point:
	
	\begin{center}
		\includegraphics[scale=0.45]{FN7.pdf} 
	\end{center}
	
	\begin{align*}
		\left(\frac{\partial g_6}{\partial x}\right)(x,y,z)&= \textcolor{red}{h_{14}h_{46}} \frac{\partial g_1}{\partial x}(x,y,z)+  \textcolor{red}{\left(h_{24}h_{46} + h_{25}h_{56}\right)}  \frac{\partial g_2}{\partial x}(x,y,z)+  \textcolor{red}{h_{25}h_{56}}  \frac{\partial g_3}{\partial x}(x,y,z)
	\end{align*}
	
	One more propagation and we are done
\end{frame}


\begin{frame}
	In the final propagation we compute the derivatives of $g_1$, $g_2$ and $g_3$, but in this case there is no more chain rule to apply.
	\begin{align*}
		\left(\frac{\partial g_1}{\partial x}\right)(x) = \left(\frac{\partial f_1}{\partial x}\right)(x) = \textcolor{red}{h_{x1}}\\
		\left(\frac{\partial g_2}{\partial x}\right)(x) = \left(\frac{\partial f_2}{\partial x}\right)(x)= \textcolor{red}{h_{x2}}\\
	\end{align*}
	
	\begin{center}
		\includegraphics[scale=0.45]{FN8.pdf} 
	\end{center}
	
	
\end{frame}


\begin{frame}
	\begin{center}
		\includegraphics[scale=0.45]{FN8.pdf} 
	\end{center}
	
	Up to this point we have
	
	\begin{align*}
		\left(\frac{\partial g_6}{\partial x}\right)(x,y,z)&= \textcolor{red}{h_{14}h_{46}} \frac{\partial g_1}{\partial x}(x,y,z)+  \textcolor{red}{\left(h_{24}h_{46} + h_{25}h_{56}\right)}  \frac{\partial g_2}{\partial x}(x,y,z)+  \textcolor{red}{h_{25}h_{56}}  \frac{\partial g_3}{\partial x}(x,y,z)\\&= \textcolor{red}{h_{x1}h_{14}h_{46}} + \textcolor{red}{h_{x2}h_{24}h_{46}} + \textcolor{red}{h_{x2}h_{25}h_{56}}    
	\end{align*}
	
	In practice, we do not need to pass $h_{vw}$ but rather the cumulative product.
\end{frame}

\begin{frame}{Back Propagation Algorithm}
	\Large
	\fancythm{Back Propagation Algorithm}
	{
		\begin{enumerate}
			\item Compute the forward evaluations $g_v$ for each non-initial vertex $V$.
			\item For each vertex $v$ different than the final vertex $v$, set the variable $c_v = -\infty$ and for each $w$ that can be reached from $v$ by a direct edge set $c_{vw}= -\infty$ 
			\item Set each vertex of the graph as \textbf{not-ready}, and set the final vertex as \textbf{ready}
			\item Repeat until al vertices are \textbf{done}
			\begin{enumerate}
				\item[i.] For all $w \in V$ such that $w$ is \textbf{ready}, and all $v \in V$ such that $v \to w$, then $w$ propagates backward to $v$ the value $c_{vw}=c_{w} \times h_{vw}$ with $h_{vw} = (\partial_v f_w)(g_v)$ where $\partial_v f_w$ means the derivative with respect to the coordinate where $v$ is included in $f_w$.
				\item[ii.] Set all \textbf{ready} vertices as \textbf{done}. 
				\item[iii.] For all vertices $v \in V$ that are \textbf{non-ready}, and check if they have receive all values $c_{vw}$. In such a case, set $c_v = \sum_{w: v\to w} c_{vw}$ and $v$ is set to be \textbf{ready}
			\end{enumerate}
		\end{enumerate}
	}
\end{frame}

\begin{frame}{Back Propagation Algorithm}
	In practice, for implementing the gradient descent algorithm we need. 
	
	\begin{enumerate}
		\item Gradient of each function presented in a vertex of the network
		\item Each vertex has a label variable indicating that the vertex is \textbf{not-ready}, \textbf{ready} or \textbf{done}
		\item Each vertex has a list $(c_{vw}: w\in V, v\to w)$, and a value $c_v$. Set $c_f = 1$ for the final vertex.
	\end{enumerate}
\end{frame}

\begin{frame}{Back Propagation Algorithm: final comments}
	
	\begin{witemize}
		\item Back-propagation is a very good way to keep track of partial derivatives and it can be used in a lot of algorithms that require derivatives, not only in Gradient Descent algorithms.
		
		\item It automatises the computation of derivatives that sometimes are painful to do by hand
		
		\item Automatic differentiation is a very important area of research.
		
		\item In some cases we don't want to implement the algorithm but use it only to \textit{\textcolor{blue}{get ideas}} of how to compute gradients. If we have a well defined model BP can help us to be careful and clean with our ideas.
		
		\item In some other cases BP is unavoidable.
	\end{witemize}
	
\end{frame}


\end{document}